\documentclass[10pt,a4paper]{amsart}
\usepackage{amsmath,amssymb,amsthm,graphicx}
% Retained preamble for the Stacks Project chapter divisors.tex (Divisors).
% The chapter file pinned-original-source.tex carries only \input{preamble} on its first line, so no
% preamble text lives inside the chapter itself.  The Stacks Project uses one single shared preamble for
% every chapter; the snapshot stages it verbatim in sources/stacks-chow/preamble.tex.  The block below is
% copied mechanically, line for line, from that file's lines 14-18 (the Xy-pic setup) and lines 154-217
% (the native theorem declarations together with the mathematical macros).  Nothing was added, removed or
% redefined.  The wrapper class and the amsmath/amssymb/amsthm/graphicx base packages are supplied by the
% assembler, so the class/\externaldocument part of the shared preamble is not repeated here.
% For commutative diagrams we use Xy-pic
\usepackage[all]{xy}

% We use 2cell for 2-commutative diagrams.
\xyoption{2cell}
% Theorem environments.
%
\theoremstyle{plain}
\newtheorem{theorem}[subsection]{Theorem}
\newtheorem{proposition}[subsection]{Proposition}
\newtheorem{lemma}[subsection]{Lemma}

\theoremstyle{definition}
\newtheorem{definition}[subsection]{Definition}
\newtheorem{example}[subsection]{Example}
\newtheorem{exercise}[subsection]{Exercise}
\newtheorem{situation}[subsection]{Situation}

\theoremstyle{remark}
\newtheorem{remark}[subsection]{Remark}
\newtheorem{remarks}[subsection]{Remarks}

\numberwithin{equation}{subsection}

% Macros
%
\def\lim{\mathop{\mathrm{lim}}\nolimits}
\def\colim{\mathop{\mathrm{colim}}\nolimits}
\def\Spec{\mathop{\mathrm{Spec}}}
\def\Hom{\mathop{\mathrm{Hom}}\nolimits}
\def\Ext{\mathop{\mathrm{Ext}}\nolimits}
\def\SheafHom{\mathop{\mathcal{H}\!\mathit{om}}\nolimits}
\def\SheafExt{\mathop{\mathcal{E}\!\mathit{xt}}\nolimits}
\def\Sch{\mathit{Sch}}
\def\Mor{\mathop{\mathrm{Mor}}\nolimits}
\def\Ob{\mathop{\mathrm{Ob}}\nolimits}
\def\Sh{\mathop{\mathit{Sh}}\nolimits}
\def\NL{\mathop{N\!L}\nolimits}
\def\CH{\mathop{\mathrm{CH}}\nolimits}
\def\proetale{{pro\text{-}\acute{e}tale}}
\def\etale{{\acute{e}tale}}
\def\QCoh{\mathit{QCoh}}
\def\Ker{\mathop{\mathrm{Ker}}}
\def\Im{\mathop{\mathrm{Im}}}
\def\Coker{\mathop{\mathrm{Coker}}}
\def\Coim{\mathop{\mathrm{Coim}}}

% Boxtimes
%
\DeclareMathSymbol{\boxtimes}{\mathbin}{AMSa}{"02}

%
% Macros for moduli stacks/spaces
%
\def\QCohstack{\mathcal{QC}\!\mathit{oh}}
\def\Cohstack{\mathcal{C}\!\mathit{oh}}
\def\Spacesstack{\mathcal{S}\!\mathit{paces}}
\def\Quotfunctor{\mathrm{Quot}}
\def\Hilbfunctor{\mathrm{Hilb}}
\def\Curvesstack{\mathcal{C}\!\mathit{urves}}
\def\Polarizedstack{\mathcal{P}\!\mathit{olarized}}
\def\Complexesstack{\mathcal{C}\!\mathit{omplexes}}
% \Pic is the operator that assigns to X its picard group, usage \Pic(X)
% \Picardstack_{X/B} denotes the Picard stack of X over B
% \Picardfunctor_{X/B} denotes the Picard functor of X over B
\def\Pic{\mathop{\mathrm{Pic}}\nolimits}
\def\Picardstack{\mathcal{P}\!\mathit{ic}}
\def\Picardfunctor{\mathrm{Pic}}
\def\Deformationcategory{\mathcal{D}\!\mathit{ef}}

% --- appended 2026-10-12 (Chief): the project own definitions of the reference/slogan/history
% display environments, copied mechanically from the shared Stacks preamble lines 8-12
% (\usepackage{verbatim} plus the three \newenvironment lines, which the chapter preamble keeps
% verbatim from the project preamble). Nothing added, removed or redefined: these environments
% carry no mathematical content and the project itself typesets them as comments. Without these
% three lines a counted statement whose own text uses them fails with Environment slogan undefined.
\usepackage{verbatim}
\newenvironment{reference}{\comment}{\endcomment}
\newenvironment{slogan}{\comment}{\endcomment}
\newenvironment{history}{\comment}{\endcomment}
\usepackage{cleveref}
\crefname{originallemma}{original lemma}{original lemmas}
\Crefname{originallemma}{Original lemma}{Original lemmas}
\makeatletter
\newlabel{more-algebra-lemma-fitting-ideal-basics}{{more-algebra-lemma-fitting-ideal-basics}{source}{}{sourcepointer-1}{}}
\newlabel{more-algebra-lemma-fitting-ideal-basics@cref}{{[originallemma][0][]more-algebra-lemma-fitting-ideal-basics}{[1][0][]source}{}{}{}}
\makeatother
\crefname{originaldefinition}{original definition}{original definitions}
\Crefname{originaldefinition}{Original definition}{Original definitions}
\makeatletter
\newlabel{morphisms-definition-scheme-theoretic-support}{{07U1}{source}{}{sourcepointer-2}{}}
\newlabel{morphisms-definition-scheme-theoretic-support@cref}{{[originaldefinition][0][]07U1}{[1][0][]source}{}{}{}}
\makeatother
\crefname{originallemma}{original lemma}{original lemmas}
\Crefname{originallemma}{Original lemma}{Original lemmas}
\makeatletter
\newlabel{morphisms-lemma-i-star-equivalence}{{04CJ}{source}{}{sourcepointer-3}{}}
\newlabel{morphisms-lemma-i-star-equivalence@cref}{{[originallemma][0][]04CJ}{[1][0][]source}{}{}{}}
\makeatother
\crefname{originallemma}{original lemma}{original lemmas}
\Crefname{originallemma}{Original lemma}{Original lemmas}
\makeatletter
\newlabel{morphisms-lemma-scheme-theoretic-support}{{07U0}{source}{}{sourcepointer-4}{}}
\newlabel{morphisms-lemma-scheme-theoretic-support@cref}{{[originallemma][0][]07U0}{[1][0][]source}{}{}{}}
\makeatother
\title{The closed subscheme cut out by the zeroth Fitting ideal agrees with the scheme theoretic support and carries the module}
\author{The Stacks Project Authors}
\date{stacks-project:divisors.tex:local-snapshot-20261012; GNU FDL 1.2 or later}
\begin{document}
\maketitle
\noindent\textit{Editorial note.} All delivered mathematical text is the author's own, copied from the pinned Stacks Project chapter file pinned-original-source.tex; no statement or proof step was rewritten, reordered or removed.\par\medskip
\noindent\textit{Editorial note.} Following the multi-claim convention for this source, source\_url is the source-level URL https://stacks.math.columbia.edu/ and the source label lemma-on-subscheme-cut-out-by-Fit-0 is carried in claim\_key and here rather than in a per-claim tag URL.\par\medskip
\noindent\textit{Editorial note.} The preamble is the chapter's shared Stacks Project preamble, retained as sources/stacks-divisors/retained-preamble.tex: a line-for-line excerpt of the shared Stacks preamble staged with this snapshot. No macro was invented or redefined, and this card adds no preamble text of its own.\par\medskip
\noindent\textit{Editorial note.} Statement span 1309-1326 and proof span 1328-1349 come from the halfA pool card operations/pool-cards/stacks-divisors-halfA.json and were re-read in the pinned chapter file before this card was written.\par\medskip
\noindent\textit{Editorial note.} The proof offers two ways of producing the module of part (3) with an 'either ... or ...' wording; both routes are the author's and are preserved.\par\medskip
\noindent\textit{Editorial note.} 4 cross-references occur inside the counted spans; they are listed in external\_references as navigation pointers only.\par\medskip
\noindent\textit{Allowed standard prerequisites (editorial).}\begin{itemize}
\item The scheme theoretic support of a finite type quasi-coherent module and its local description by the annihilator, as defined and established in the morphisms chapter. Cited at proof lines 1329-1332 and 1338-1340. Identifies the scheme theoretic support and its support as a closed subset, which are two of the claims of the statement.
\item The elementary properties of Fitting ideals, in particular that the zeroth Fitting ideal annihilates the module and that the closed subscheme it cuts out has the same underlying set as the support. Cited at proof lines 1332-1334 and 1340-1342. Gives the scheme theoretic inclusion of the support in the Fitting subscheme and the equality of the two underlying closed subsets.
\item The equivalence between quasi-coherent modules on a closed subscheme and quasi-coherent modules annihilated by the ideal sheaf, cited from the morphisms chapter. Cited at proof lines 1344-1347. Produces the finite type module on the Fitting subscheme whose pushforward is the given module.
\end{itemize}\medskip
\ifdefined\hypertarget\hypertarget{sourcepointer-1}{}\fi
\noindent\textit{Source reference.} Original lemma more-algebra-lemma-fitting-ideal-basics: sources/stacks-divisors/pinned-original-source.tex, source lines 1266--1266. Same-chapter reference used in the retained author text at source line(s) 1266, 1294, 1306, 1334, 1340, 1348. It is supplied only as a navigation pointer so that the retained text compiles; the counted item delivers the author's own statement and the author's complete proof, and no part of the referenced item is copied in or substituted for the counted argument.\par\medskip
\ifdefined\hypertarget\hypertarget{sourcepointer-2}{}\fi
\noindent\textit{Source reference.} Original definition 07U1: focused-seeds/AGP/sources/stacks-snapshot-20261011/tags.txt, source lines 9956--9956. Same-chapter reference used in the retained author text at source line(s) 1330. It is supplied only as a navigation pointer so that the retained text compiles; the counted item delivers the author's own statement and the author's complete proof, and no part of the referenced item is copied in or substituted for the counted argument.\par\medskip
\ifdefined\hypertarget\hypertarget{sourcepointer-3}{}\fi
\noindent\textit{Source reference.} Original lemma 04CJ: focused-seeds/AGP/sources/stacks-snapshot-20261011/tags.txt, source lines 5693--5693. Same-chapter reference used in the retained author text at source line(s) 1345, 2728, 4308. It is supplied only as a navigation pointer so that the retained text compiles; the counted item delivers the author's own statement and the author's complete proof, and no part of the referenced item is copied in or substituted for the counted argument.\par\medskip
\ifdefined\hypertarget\hypertarget{sourcepointer-4}{}\fi
\noindent\textit{Source reference.} Original lemma 07U0: focused-seeds/AGP/sources/stacks-snapshot-20261011/tags.txt, source lines 9955--9955. Same-chapter reference used in the retained author text at source line(s) 1331, 1337, 1343. It is supplied only as a navigation pointer so that the retained text compiles; the counted item delivers the author's own statement and the author's complete proof, and no part of the referenced item is copied in or substituted for the counted argument.\par\medskip
\begin{lemma}
\label{lemma-on-subscheme-cut-out-by-Fit-0}
Let $S$ be a scheme.
Let $\mathcal{F}$ be a finite type, quasi-coherent $\mathcal{O}_S$-module.
Let $Z_0 \subset S$ be the closed subscheme cut out by
$\text{Fit}_0(\mathcal{F})$.
Let $Z \subset S$ be the scheme theoretic support of $\mathcal{F}$.
Then
\begin{enumerate}
\item $Z \subset Z_0 \subset S$ as closed subschemes,
\item $Z = Z_0 = \text{Supp}(\mathcal{F})$ as closed subsets,
\item there exists a finite type, quasi-coherent $\mathcal{O}_{Z_0}$-module
$\mathcal{G}_0$ with
$$
(Z_0 \to X)_*\mathcal{G}_0 = \mathcal{F}.
$$
\end{enumerate}
\end{lemma}

\begin{proof}
Recall that $Z$ is locally cut out by the annihilator of $\mathcal{F}$, see
Morphisms, Definition \ref{morphisms-definition-scheme-theoretic-support}
(which uses Morphisms, Lemma \ref{morphisms-lemma-scheme-theoretic-support}
to define $Z$). Hence we see that $Z \subset Z_0$ scheme theoretically
by More on Algebra, Lemma
\ref{more-algebra-lemma-fitting-ideal-basics} part (6).
On the other hand we have $Z = \text{Supp}(\mathcal{F})$
set theoretically by
Morphisms, Lemma \ref{morphisms-lemma-scheme-theoretic-support}
and we have $Z_0 = Z$ set theoretically by
More on Algebra, Lemma
\ref{more-algebra-lemma-fitting-ideal-basics} part (7).
Finally, to get $\mathcal{G}_0$ as in part (3) we can either use
that we have $\mathcal{G}$ on $Z$ as in
Morphisms, Lemma \ref{morphisms-lemma-scheme-theoretic-support}
and set $\mathcal{G}_0 = (Z \to Z_0)_*\mathcal{G}$
or we can use Morphisms, Lemma \ref{morphisms-lemma-i-star-equivalence}
and the fact that $\text{Fit}_0(\mathcal{F})$ annihilates
$\mathcal{F}$ by More on Algebra, Lemma
\ref{more-algebra-lemma-fitting-ideal-basics} part (6).
\end{proof}

\end{document}
